%----------------------------------------------------------------------------
\chapter{Előzmények}
%----------------------------------------------------------------------------
Korábban említettem, hogy önálló labor tantárgy kerertein belül foglalkoztam
a felhő natív architektúra elemeivel egy konkrét webapplikáció fejlesztése révén,
amely egy weboldalra könnyen beágyazható csevegőrobotot takart. Az akkori
feladat során mély betekintést nyertem a felhőnatív applikációk világába, illetve
a nagy nyelvi modellek használatába. A korábbi és ezen projekt között sok átfedés
van a használt technológiák terén. Az egyik ilyen legfontosabb átfedő elem az API.

A csevegőrobot rendszer működésének központi elemét az API alkotja,
amely hídként szolgál az elosztott programkomponensek között.
Ez a felület felel a beérkező klienskérések fogadásáért, érvényesítéséért és a 
megfelelő háttérszolgáltatásokhoz történő irányításáért. 
Gyakorlati szempontból az API az a csomópont, amely a tényleges 
felhasználói interakciókat kezeli: ezeken a végpontokon keresztül történik például a 
csevegőrobotnak küldött üzenetek fogadása,
majd a generált válaszok visszatovábbítása a felhasználó felé. Ugyancsak ez a komponens
biztosítja a tudásbázis adminisztrációs háttérműveleteit is, ideértve az új dokumentumok
névlegesen PDF-ek (Portable Document Format) feltöltését, azok feldolgozását, valamint a 
későbbi információs lekérdezések lebonyolítását.
Emellett az adatközvetítés szabványosított JSON (JavaScript Object Notation ) formátumban,
aszinkron, másnéven nem blokkoló hívásokon keresztül történik, ami lehetővé teszi, hogy a rendszer
párhuzamosan több kérést is kiszolgáljon, ez
javíthatja a rendszer teljesítményét az olyan időigényes háttérfolyamatoknál is, 
mint a nagy méretű PDF dokumentumok feldolgozása vagy a nagy nyelvi modellel való 
kommunikáció.

Az általam megvalósított felhőoptimalizáló rendszerben az API a felhasználói felület
és az alkalmazás háttérlogikája közötti kapcsolatot 
biztosítja. Ezen keresztül lehet lekérdezni a megfigyelt alkalmazások állapotát, 
az optimalizálási javaslatokat és a korábbi műveletek eredményeit. Szintén ezen 
a felületen keresztül érkeznek a felhasználói kérések, például egy új elemzés indítása,
egy javaslat jóváhagyása vagy elutasítása, illetve egy korábbi módosítás visszaállítása.
Az API fogadja ezeket a kéréseket, ellenőrzi a megadott paramétereket,
majd átadja őket a megfelelő feldolgozó résznek.

Fontos különbség azonban, hogy ezek a feldolgozó részek nem önálló mikroszolgáltatásokként,
hanem egy közös alkalmazáson belül működnek.
A háttérlogika a megfigyelő által gyűjtött mérési adatokat használja fel,
a nagy nyelvi modelltől optimalizálási javaslatot kér, majd felhasználói jóváhagyás után, 
biztonsági ellenőrzések mellett módosítja az orkesztrátor által futtatott alkalmazás
beállításait.

Az API-hoz mindkét esetben a FastAPI-t használtam, amelyről majd később lesz bővebben szó.\\

A két alkalmazásban egy további közös vonás valamilyen adatbázis használata,
azonban ezek szerepe eltér egymástól. A csevegőrobot esetében a vektoradatbázis a
tudásbázisban történő keresést segíti. Ennek célja, hogy a felhasználó kérdéséhez
kapcsolódó információk bekerüljenek a nagy nyelvi modell számára átadott kontextusba, 
így a válasz ne kizárólag a modell korábban megszerzett ismereteire támaszkodjon. 
Ez csökkentheti a hallucináció jelenségét, amikor a modell olyan állításokat fogalmaz meg,
amelyek tévesek vagy nem támaszthatók alá a rendelkezésre álló információkkal.

A visszakeresés és a szöveggenerálás összekapcsolására szolgál az úgynevezett
RAG \cite{NEURIPS2020_6b493230} (Retrieval-Augmented Generation), vagyis visszakereséssel kiegészített szöveggenerálás.
Ennek során a csevegő rendszer először megkeresi a kérdéshez kapcsolódó információkat egy külső
tudásforrásban, majd ezekkel egészíti ki a modell bemenetét, ezzel csökkentve a hallucináció
fellépésének esélyét. 

Az optimalizáló esetében szintén elengedhetetlen a megfelelő kontextus biztosítása,
de ezt nem dokumentumok visszakeresésével olda meg, hanem a rendszer a megfigyelt
alkalmazások mérési adatait, aktuális erőforrás-beállításait és az előre meghatározott
szabályokat közvetlenül átadja modellnek egy rendszerutasítással együtt, így képezve megfelelő kontextust.
Ezek alapján a modell konkrét optimalizálási javaslatot fogalmazhat meg.
Az adatbázis itt elsősorban a javaslatok, a felhasznált mérési
adatok és a végrehajtott műveletek megőrzésére szolgál, hogy a rendszer működése később is
visszakövethető legyen.

A megfelelő kontextus mellett a javaslatok ellenőrzése is fontos. A rendszer ezért külön
ellenőrzi, hogy a modell által megadott értékek megfelelnek-e az előírt korlátoknak,
a tényleges módosítás pedig csak felhasználói jóváhagyás után történhet meg.
